\subsection{Lyapunov-Style Reasoning}

\begin{definition}[Lyapunov-style certificate]
A \textbf{Lyapunov-style certificate} is a scalar quantity that decreases along
trajectories and thereby provides evidence that the system is moving toward a
stable set.
\end{definition}

The point of Lyapunov language in this book is not to carry out rigorous proofs
for everyday routines. The point is to introduce a design ideal: can the
intervention supply a monotone indicator of progress rather than relying on
vague self-report?

\begin{example}[Monotone progress proxy]
Suppose a reader wants to establish a reading regime. A useful practical proxy
is not ``did I feel scholarly?'' but ``did I produce the first margin note
before a fixed time?'' That is not a mathematical Lyapunov proof. It is a
Lyapunov-style design move: choose a measurable quantity whose consistent
improvement supports the hypothesis that the target regime is becoming easier to
reach.
\end{example}
